\chapter{Kata Pengantar}

Buku ini memiliki satu cita-cita sederhana: sebuah buku ilmu kebumian dan
antariksa yang utuh dan berkesinambungan, mencakup segala hal sejak tahun
pertama sekolah hingga akhir program sarjana, yang ditulis untuk pembaca
di seluruh dunia, mula-mula dalam bahasa Inggris, lalu diterjemahkan di
sini ke dalam bahasa Indonesia.

Gayanya ringkas dan cermat: pelajaran dibangun dari definisi, contoh,
proposisi, metode, dan diagram, disertai penalaran yang cermat setiap kali
penalaran itu terjangkau pada tingkat yang bersangkutan, lalu diikuti
latihan yang dipilih dengan saksama. Hasil yang dinyatakan tanpa penurunan
lengkap ditandai secara tegas sebagai diterima tanpa bukti. Penyelesaian
lengkap semua latihan dikumpulkan di akhir buku. Setiap angka yang dikutip
tentang Bumi yang nyata, samudranya, atmosfernya, atau langit berasal dari
sumber acuan, dan setiap kurva yang digambar dari hasil pengukuran nyata
dihitung dari data yang telah diterbitkan.

Pelajaran ini terbagi atas lima buku. Buku pertama mencakup tahun-tahun
sekolah dasar dan sekolah menengah pertama, buku kedua ketiga tahun
sekolah menengah atas, dan tiga buku terakhir ketiga tahun universitas.
Kedua buku sekolah tidak pernah mengajarkan gagasan yang sama dua kali:
setiap gagasan diajarkan sekali, pada tahun yang pertama kali cocok
untuknya, dan bab berikutnya yang memerlukannya dibuka dengan pengingat
singkat, \emph{Yang sudah kamu ketahui}, sebelum melangkah lebih jauh.
Jika suatu gagasan termasuk fisika, biologi, atau kimia, buku-buku
pendamping dalam koleksi yang sama mengajarkannya, dan buku ini
mengingatkannya kembali sebelum memandangnya dari sudut pandang seorang
ilmuwan kebumian. Dalam buku-buku sekolah, setiap tahun menjadi satu
bagian, yang terbagi atas beberapa bab; setiap buku universitas adalah
satu tahun, yang terbagi atas beberapa bab. Semakin awal tahunnya,
semakin muda pembaca yang dituju: lebih banyak ilustrasi, langkah
pengerjaan yang lebih terperinci, dan latihan yang lebih ringan.

\section*{Cara membaca buku ini}

Pernyataan diberi nomor di dalam setiap bab dan dibedakan dengan warna:
definisi berwarna biru, teorema merah, proposisi beserta kerabatnya
jingga, dan metode --- resep singkat untuk pekerjaan yang baku ---
berwarna hijau. Latihan bertingkat dari ($\star$) untuk penerapan langsung
sampai ($\star\star\star$) untuk soal yang lebih sukar. Kotak berbingkai
memuat hal-hal yang bukan bagian dari alur penalaran: pengingat tentang
apa yang telah diajarkan bab sebelumnya, cara para ahli geologi dan
ilmuwan kebumian lainnya bekerja di lapangan atau di laboratorium, satu
halaman sejarah, angka-angka yang berubah dari tahun ke tahun (beserta
tahun berlakunya), dan nasihat keselamatan.

\section*{Keselamatan}

Bumi dipelajari dengan mengamatinya: langit, cuaca, bentang alam, batu
pada sebuah bangunan, tanah di sebuah kebun. Hanya itulah yang pernah
diminta buku ini dari pembacanya. Jangan pernah memandang Matahari, baik
dengan maupun tanpa alat: Matahari dapat membutakan. Tambang batu,
tebing, gua, sungai yang sedang banjir, dan gunung api adalah tempat bagi
para ahli yang terlatih dan berperlengkapan, dan buku ini menguraikan
pekerjaan mereka tanpa menyuruh siapa pun pergi ke sana.

\section*{Berkontribusi}

Kode sumber buku ini tersimpan di sebuah repositori git publik:
\begin{center}
\url{https://github.com/bvirrion/one-earth-and-space-sciences-book}
\end{center}
Kontribusi --- bab baru, koreksi, penjelasan yang lebih baik, latihan
tambahan --- sangat kami sambut; lihat \texttt{CONTRIBUTING.md} di akar
repositori.
